\section{Conclusion}

This paper studies return prediction from a matched 20-day OHLCV information set. The evidence does not support a simple view that chart-image CNNs universally dominate numerical models. Strong nonlinear tabular learners can match or exceed CNN performance, and a 1D-CNN trained directly on the OHLCV sequence is stronger still. All Sharpe ratios are computed from arithmetic mean period returns and period volatility, annualized using 52 weekly periods; compounded annual returns are reported separately.

However, this does not make image models redundant. CNN, MLP, HistGB, and temporal numerical scores are only moderately correlated, retain incremental predictive content after controls, and improve performance when combined. We use two non-interchangeable residual estimands: a three-model CNN--MLP-4--HistGB decomposition and a pairwise CNN--1D-CNN temporal decomposition. The 1D-CNN absorbs a large share of the CNN-versus-flattened-model residual, but the CNN-specific component conditional on 1D-CNN remains positive and the CNN plus 1D-CNN ensemble outperforms either standalone model. In that pairwise temporal decomposition, the residual chart-CNN signal is short-horizon and locally path/volume sensitive, whereas the residual 1D-CNN signal is persistent and pressure-sensitive. These results identify complementary transformations of the same price history, rather than a universal advantage of image representation.

The main implication is methodological and economic. Model comparisons should not stop at asking which learner has the highest Sharpe ratio. When models share the same information set, their disagreements can reveal how different function classes parse price history into complementary trading signals. The liquidity audit shows that predictability is stronger among smaller stocks, but remains economically meaningful in a large-and-liquid universe; simple turnover-proportional costs nevertheless materially reduce net performance. Thus the results establish complementary OOS predictability under the stated portfolio protocol, not unlimited or frictionless arbitrage capacity. Further robustness should examine stronger temporal numerical architectures, explicit execution constraints, price impact, and external markets.
